% Part: applied-modal-logics
% Chapter: epistemic-logic
% Section: public-announcement-logic-semantics

\documentclass[../../../include/open-logic-section]{subfiles}

\begin{document}

\olfileid{aml}{el}{psm}

\olsection{معناشناسی منطق اعلان عمومی}

مدل‌های رابطه‌ایِ منطق‌های اعلان عمومی همان مدل‌های منطق معرفتی‌اند.
اما معناشناسی عملگر اعلان عمومی امری تازه است.

\begin{defn}\ollabel{defn:mmodels} \emph{صدق !!a{formula}~$!A$
  در~$w$} در مدل~$\mModel M = \tuple{W, R, V}$، که با
  $\mSat{M}{!A}[w]$ نمادگذاری می‌شود، به‌طور استقرایی چنین تعریف می‌شود:
  \begin{enumerate}
  \tagitem{prvFalse}{\indcase{!A}{\lfalse}{$\mSat{M}{\lfalse}[w]$ هرگز برقرار نیست.}}{}
  \tagitem{prvTrue}{\indcase{!A}{\ltrue}{$\mSat{M}{\ltrue}[w]$ همواره برقرار است.}}{}
  \item $\mSat{M}{p}[w]$ اگر و تنها اگر~$w \in V(p)$.
  \tagitem{prvNot}{\indcase{!A}{\lnot !B}{$\mSat{M}{\indfrm}[w]$ اگر و تنها اگر
    $\mSat/{M}{!B}[w]$.}}{}
  \tagitem{prvAnd}{\indcase{!A}{(!B \land !C)}{$\mSat{M}{\indfrm}[w]$ اگر و تنها اگر
    $\mSat{M}{!B}[w]$ و~$\mSat{M}{!C}[w]$.}}{}
  \tagitem{prvOr}{\indcase{!A}{(!B \lor !C)}{$\mSat{M}{\indfrm}[w]$ اگر و تنها اگر
    $\mSat{M}{!B}[w]$ یا~$\mSat{M}{!C}[w]$} (یا هر دو).}{}
  \tagitem{prvIf}{\indcase{!A}{(!B \lif !C)}{$\mSat{M}{\indfrm}[w]$ اگر و تنها اگر
    $\mSat/{M}{!B}[w]$ یا~$\mSat{M}{!C}[w]$.}}{}
  \tagitem{prvIff}{\indcase{!A}{(!B \liff !C)}{$\mSat{M}{\indfrm}[w]$ اگر و تنها اگر
    یا هر دو رابطهٔ~$\mSat{M}{!B}[w]$ و~$\mSat{M}{!C}[w]$ برقرار باشند، یا
    هیچ‌یک از~$\mSat{M}{!B}[w]$ و~$\mSat{M}{!C}[w]$ برقرار نباشند.}}{}
  \item\ollabel{defn:sub:mmodels-box}
    \indcase{!A}{\Knows_a !B}{$\mSat{M}{\indfrm}[w]$ اگر و تنها اگر
    برای هر~$w' \in W$ که~$R_a ww'$، رابطهٔ~$\mSat{M}{!B}[w']$ برقرار باشد.}
  \item\ollabel{defn:sub:mmodels-pal}
    \indcase{!A}{[!B] !C}{$\mSat{M}{\indfrm}[w]$ اگر و تنها اگر
    $\mSat{M}{!B}[w]$ مستلزم~$\mSat{M \mid !B}{!C}[w]$ باشد.}

  در اینجا~$\mModel M \mid !B = \tuple{W', R', V'}$ چنین تعریف می‌شود:
  \begin{enumerate}
    \item $W' = \Setabs{ u \in W}{\mSat{M}{!B}[u]}$. پس جهان‌های
      $\mModel M \mid !B$ همان جهان‌های~$\mModel M$ هستند که~$!B$
      در آن‌ها برقرار است.
    \item $R'_a = R_a \cap (W' \times W')$. رابطهٔ دسترسی‌پذیریِ هر
      عامل صرفاً به جهان‌های باقی‌مانده در~$W'$ محدود می‌شود.
    \item $V'(p) = \Setabs{ u \in W'}{u \in V(p) }$. به همین ترتیب،
      ارزش‌گذاری‌های گزاره‌ای در جهان‌ها یکسان می‌مانند؛ این امر بیانگر
      آن است که رویدادهای اطلاعاتی ارزش صدق متغیرهای گزاره‌ای را تغییر
      نمی‌دهند.
  \end{enumerate}
  
  \end{enumerate}
\end{defn}

پس ویژگی متمایز منطق‌های اعلان عمومی آن است که گاهی صدق
!!a{formula} در~$\mModel M$ را تنها با ارجاع به مدلی غیر از خودِ
$\mModel M$ می‌توان تعیین کرد.

همچنین توجه کنید که معناشناسی ما عملگر اعلان را عملگری از نوع~$\Box$
می‌گیرد. بنابراین، اگر نتوان !!a{formula}~$!A$ را در جهانی صادقانه
اعلان کرد، آنگاه~$[!A]!B$ در آنجا بدیهی‌وار برقرار خواهد بود، درست
همان‌گونه که همهٔ !!{formula}های~$\Box$ در نقاط پایانی برقرارند.

\begin{figure}
  \begin{center}
    \begin{tikzpicture}[modal]
      \node[world] (w1) [label={below left:$p, \lnot q$}]{$w_1$};
      \node[world] (w2) [left=of w1, label={left:$\lnot p, \lnot q$}]{$w_2$};
      \node[world] (w3) [above=of w1, label={right:$p, q$}] {$w_3$};
      \draw[<->] (w1) to [swap] node {$b$} (w2);
      \draw[->,loop below] (w1) to node {$a, b$} (w1);
      \draw[<->] (w1) to [swap] node {$a$} (w3);
      \draw[->,loop above] (w2) to node {$a, b$} (w2) ;
      \draw[->,loop above] (w3) to node {$a, b$} (w3) ;
      
      \node[world](v1)[right of=w1, xshift=1.25in, label={right:$p, \lnot q$}]{$w'_1$};
      \node[world](v3)[above of=v1,label={right:$p,q$}]{$w'_3$};
      \draw[<->] (v1) to node {$a$} (v3);
      \draw[->,loop below] (v1) to node {$a,b$} (v1);
      \draw[->,loop above] (v3) to node {$a,b$} (v3) ;
      
      \node(m)[below=of w1]{$\mModel M$};
      \node(m')[below=of v1]{$\mModel M \mid p$};

      \draw[-,dotted] (w1) to node {\footnotesize اعلانِ~$p$}(v1) ;
     
    \end{tikzpicture}
  \end{center}
  \caption{پیش و پس از اعلان عمومیِ~$p$.}
  \ollabel{fig:announcement-example}
\end{figure}

می‌توان اعلان عمومیِ !!a{formula} را کوچک‌کردن یک مدل، یا محدودکردن
آن به جهان‌هایی دانست که !!{formula} در آن‌ها صادق بوده است.
\olref{fig:announcement-example} نمونه‌ای از آثار اعلان عمومیِ~$p$ را
نشان می‌دهد. نکتهٔ شایان توجه دربارهٔ آن مدل این است که در نتیجهٔ
اعلان، عامل~$b$ درمی‌یابد که~$p$، اما عامل~$a$ چنین دانشی را تازه به
دست نمی‌آورد (زیرا~$a$ از پیش می‌دانست که~$p$ صادق است).

به‌طور صوری، $\mSat{M}{\lnot \Knows_b p}[w_1]$، اما
$\mSat{M \mid p}{\Knows_b p}[w_1]$. از این امر نتیجه می‌شود که
$\mSat{M}{[p] \Knows_b p}[w_1]$. اما دربارهٔ نتیجهٔ اعلان می‌توانیم
ادعاهایی حتی قوی‌تر مطرح کنیم. در واقع،
$\mSat{M}{[p]\CKnows_{\{a,b\}} p}[w_1]$. به بیان دیگر، پس از اعلانِ
$p$، آن گزاره به \emph{دانش مشترک} بدل می‌شود.

بااین‌حال، شاید بپرسیم آیا این نتیجه در حالت کلی برقرار است و آیا اعلان
صادقانهٔ~$!A$ \emph{همواره} موجب می‌شود~$!A$ به دانش مشترک بدل گردد.
شاید شگفت‌آور باشد که پاسخ منفی است. در واقع، می‌توان !!{formula}هایی
را صادقانه اعلان کرد که پس از اعلان دیگر صادق نباشند. برای نمونه، آثار
اعلانِ~$p \land \lnot \Knows_b p$ در~$w_1$ از
\olref{fig:announcement-example} را در نظر بگیرید. مدل
$\mModel M \mid p$ جهان‌های~$w_1$ و~$w_3$ را نگه می‌دارد، اما
$\mModel M \mid (p \land \lnot \Knows_b p)$ تنها جهان~$w_1$ را نگه
می‌دارد. در این مدل تک‌جهانی،
$\mSat{M \mid (p \land \lnot \Knows_b p)}{\Knows_b p}[w_1]$؛ بنابراین
$\mSat/{M \mid (p \land \lnot \Knows_b p)}{p \land \lnot \Knows_b p}[w_1]$
و در نتیجه
$\mSat/{M}{[p \land \lnot \Knows_b p](p \land \lnot \Knows_b p)}[w_1]$.
پس این !!a{formula} پس از اعلان خود کاذب می‌شود.

\end{document}
